% \iffalse meta-comment
%
% hit-thesis-abstract.dtx 是中英文摘要页
%
% \fi
%
% \section{中英文摘要页}
%
%    \begin{macrocode}
%<*thesis-abstract>
%    \end{macrocode}
%
% 生成参考文献和关键字。
%    \begin{macrocode}
\box_new:N \hit@keywords@box
\NewDocumentCommand \hit@put@keywords { m m }
  {
    \group_begin:
      \setbox \hit@keywords@box = \hbox {#1}
      \noindent \hangindent \wd \hit@keywords@box \hangafter 1
      \box \hit@keywords@box #2 \par
    \group_end:
  }

%    \end{macrocode}
%
% \changes{v3.0.0}{2020/05/21}{英文版目录之前有page行}
% \changes{v3.0.0}{2020/05/21}{博后关键字加粗}
% \changes{v3.2a}{2026/08/07}{\cs{hit@make@abstract} 迁到 expl3。写进目录的
%   \texttt{\string~} 换成 \cs{nobreakspace}\{\}，因为 expl3 里 \texttt{\string~}
%   是普通空格}
%    \begin{macrocode}
\cs_new_protected:Npn \hit@make@abstract
  {
    \hit@clear@page
    \hit@layout@push@part:n { frontmatter }
    \hit@format@apply@heading:nn { chapter } { line-spacing = 1.25 }
    \bool_if:NT \g__hit_en_bool
      {
        \addtocontents { toc } { \nobreakspace { } \hfill \textbf { \hit@term@table@of@contents@page@en } \par }
        \addtocontents { \ext@table }
          { { \noindent \sihao \textbf { \hit@term@table@of@contents@table@en } \nobreakspace { } \hfill \textbf { \hit@term@table@of@contents@page@en } } \par }
        \addtocontents { \ext@figure }
          { { \noindent \sihao \textbf { \hit@term@table@of@contents@figure@en } \nobreakspace { } \hfill \textbf { \hit@term@table@of@contents@page@en } } \par }
      }
    \hit@appendix@chapter * [ \hit@term@abstract@heading@zh ]
      { \hit@spread@by:nn { \hit@term@abstract@heading@spread } { \hit@term@abstract@heading@zh } }
      [ \hit@term@abstract@in@english@toc@zh ]
    \pagestyle { hit@headings }
    \pagenumbering { Roman }
    \hit@abstract@zh
    % 正文与关键词之间按范例是一个空段(单独回车一行),\hitenter 模拟之
    % (1-8 页版 p4,36.8 = 19.4 + 17.25)
    \hitenter
    \hit@put@keywords
      { \bool_if:NT \g__hit_postdoc_bool { \bfseries } \heiti \hit@term@abstract@keywords@zh }
      { \hit@info@keywords@zh }
    \hit@clear@page
    % 英文摘要标题是纯西文行。Word 按行内实际字体算行高，所以这一行走西文那一档
    % ——目标上声明的 heiti 只在中文标题上成立
    \hit@format@apply@heading:nn { chapter }
      { font-zh = none , line-spacing = 1.25 }
    \hit@appendix@chapter * { \hit@term@abstract@heading@en } [ \hit@term@abstract@in@english@toc@en ]
    \hit@format@apply@heading:nn { chapter } { line-spacing = 1.25 }
    % 英文摘要正文按范例排:Word 的 1.25 倍自动行距按行内字体自然行高算,
    % 纯西文段 TNR 12bp 实测行距 17.25;字符网格的行末余量 \rightskip 只属于
    % 汉字网格行,西文段排满真版心(例右缘 510.4);汉字行撑竿会把行距顶回
    % 网格,组内停用;范例导出勾了不断字,断词罚分同关
    \group_begin:
      % 按 abstract-en 目标铺版面(行距/版宽/首行基线/断词,见 hit-format)
      \hit@format@apply@par:n { abstract-en }
      \hit@abstract@en
      \hitenter
      \hit@put@keywords { \textbf { \hit@term@abstract@keywords@en \enskip } } { \hit@info@keywords@en }
    \group_end:
%    \end{macrocode}
%
% \changes{v3.2a}{2026/08/22}{摘要排完清一次缩写的「用过」状态}
% 摘要里出现过的缩写，正文第一次再出现时仍该全写一遍——读者未必先读摘要，
% 而且规范说的「文中第一次出现」指的是正文。清不清由
% \option{abbr}/\option{reset-after-abstract} 定，默认清。
%    \begin{macrocode}
    \bool_if:NT \g__hit_abbr_reset_bool
      { \cs_if_exist:NT \glsresetall { \glsresetall } }
    \hit@layout@pop@part:
    \hit@format@apply@heading:n { chapter }
  }
%    \end{macrocode}
%
% \changes{v3.2a}{2026/08/10}{\env{denotation} 收成 \cs{item} 列表，
%   原先自己写表格的写法移到 \env{denotation*}}
% 符号表。
%
% \changes{v3.2a}{2026/08/22}{符号表提成 \env{symbols}，\env{denotation}
%   退回 v3.1e 原样只供老文档}
% 两个环境分工：\env{symbols} 是新写法，\env{denotation} 是 v3.1e 的老写法。
%
% \env{symbols} 本身就是一个 \pkg{enumitem} 的 \texttt{description} 列表，
% 一行一个符号 \cs{item}\oarg{符号}\meta{说明}。可选参数收键值：
% \option{labelwidth} 是第一列宽度（默认 2.5\,cm），\option{subheading} 给这一段
% 加个小标题——只在一页里排两段（符号加缩略语）时才用得着。不含等号的裸值当
% 宽度，\cs{begin}\marg{symbols}\oarg{3cm} 照样认。
%
% 规范里允许把符号分成几张表、各有各的列，列表排不出来，所以还有
% \env{symbols*}：标题和表号照旧，正文由用户自己写。
%
% \env{denotation} 不带列表机制，也没有星号版：v3.1e 的它就是「排标题、重置表号、
% 正文随你写」，老文档里的内容是一张手写的 \env{table}。既然 \env{symbols} 全盘
% 接手了新写法，它退回原样即可，不必现代化，也不发废弃提示。
%    \begin{macrocode}
\RequirePackage{enumitem}
\newlist { hit@symbols } { description } { 1 }
\setlist [ hit@symbols ]
  {
    nosep ,
    font = \normalfont ,
    align = left ,
    leftmargin = ! ,% 下面三个长度之和
    labelindent = 0pt ,
    labelwidth = 2.5cm ,
    labelsep* = 0.5cm ,
    itemindent = 0pt ,
  }
\tl_new:N \l__hit_symbols_width_tl
\tl_new:N \l__hit_symbols_subheading_tl
\keys_define:nn { hit / symbols / opt }
  {
    labelwidth .tl_set:N = \l__hit_symbols_width_tl ,
    subheading .tl_set:N = \l__hit_symbols_subheading_tl ,
  }
\cs_new_protected:Npn \__hit_symbols_opt:n #1
  {
    \tl_set:Nn \l__hit_symbols_width_tl { 2.5cm }
    \tl_clear:N \l__hit_symbols_subheading_tl
    \tl_if_in:nnTF {#1} { = }
      { \hit@keys@set:nn { hit / symbols / opt } {#1} }
      { \tl_if_blank:nF {#1} { \tl_set:Nn \l__hit_symbols_width_tl {#1} } }
  }
%    \end{macrocode}
%
% 每个部件各排各的章标题。只有合并成一页的那个部件例外：它自己排一个合并
% 标题，然后把闸门顶上，里头的 \env{symbols} 与 \cs{listofabbreviations*}
% 就都不再排章标题了。
%    \begin{macrocode}
\bool_new:N \g__hit_frontlist_open_bool
\tl_new:N \g__hit_frontlist_pending_tl
\cs_new_protected:Npn \hit@frontlist@open:nn #1#2
  {
    \bool_if:NF \g__hit_frontlist_open_bool
      {
        \hit@clear@page
%    \end{macrocode}
%
% 英文版排英文标题，中文版排中文标题、英文那个进英文目录。跟结论、致谢那几处
% 一样按语种分支——常量的 \texttt{-zh}／\texttt{-en} 后缀说的是「这个部件的
% 中文名与英文名」，不是「中文版用前者、英文版用后者」。
%    \begin{macrocode}
        \bool_if:NTF \g__hit_en_bool
          { \hit@appendix@chapter * {#2} [#2] }
          { \hit@appendix@chapter * {#1} [#2] }
      }
%    \end{macrocode}
%
% 段小标题:自己写的优先,没写就取合并部件预置的那一个,取完清掉。
%    \begin{macrocode}
    \tl_if_empty:NT \l__hit_symbols_subheading_tl
      {
        \tl_set_eq:NN \l__hit_symbols_subheading_tl
          \g__hit_frontlist_pending_tl
        \tl_gclear:N \g__hit_frontlist_pending_tl
      }
    \tl_if_empty:NF \l__hit_symbols_subheading_tl
      { \hit@frontlist@subheading:V \l__hit_symbols_subheading_tl }
  }
\cs_new_protected:Npn \hit@frontlist@subheading:n #1
  { \par \addvspace { \baselineskip } \noindent { \bfseries #1 } \par \nobreak }
\cs_generate_variant:Nn \hit@frontlist@subheading:n { V }
\cs_new_protected:Npn \__hit_symbols_open:
  {
    \hit@frontlist@open:nn
      { \hit@term@list@of@symbols@heading@zh } { \hit@term@list@of@symbols@heading@en }
    \setcounter { table } { 0 }
    \cs_set:Npn \thetable { \arabic { table } }% 表编号为 1
  }
\cs_new_protected:Npn \__hit_symbols_close:
  {
    \cs_set:Npn \thetable { \arabic { chapter } - \arabic { table } }% 表编号为 7-1
    \setcounter { table } { 0 }
  }
\NewDocumentEnvironment { symbols } { O{} }
  {
    \__hit_symbols_opt:n {#1}
    \__hit_symbols_open:
    \begin { hit@symbols } [ labelwidth = \l__hit_symbols_width_tl ]
  }
  { \end { hit@symbols } \__hit_symbols_close: }
\NewDocumentEnvironment { symbols* } { O{} }
  { \__hit_symbols_opt:n {#1} \__hit_symbols_open: }
  { \__hit_symbols_close: }
%    \end{macrocode}
%
% v3.1e 的 \env{denotation}：排标题、把表号改成不带章号的形式，正文随用户写。
%    \begin{macrocode}
\NewDocumentEnvironment { denotation } { O{2.5cm} }
  {
    \hit@clear@page
    \hit@appendix@chapter *
      { \hit@term@list@of@symbols@heading@zh } [ \hit@term@list@of@symbols@heading@en ]
    \setcounter { table } { 0 }
    \cs_set:Npn \thetable { \arabic { table } }
  }
  {
    \cs_set:Npn \thetable { \arabic { chapter } - \arabic { table } }
    \setcounter { table } { 0 }
  }
%    \end{macrocode}
%
% 带星号的不开列表，正文完全交给用户。可选参数收下不用，这样 v3.1e 里
% \cs{begin}\marg{denotation}\oarg{2.5cm} 的写法加个星号就能直接跑。
%    \begin{macrocode}
\NewDocumentEnvironment { denotation* } { O{2.5cm} }
  { \__hit_denotation_open: }
  { \__hit_denotation_close: }
%    \end{macrocode}
%
%    \begin{macrocode}
%</thesis-abstract>
%    \end{macrocode}
%
% \Finale
